On se place dans le plan muni d'un repère orthonormé.

On considère les points $A(-1\ ;\ 5)$, $B(3\ ;\ 5)$ et $C(4\ ;\ 0)$.

\begin{enumerate}
\item \begin{enumerate}
\item On calcule les coordonnées des vecteurs $\overrightarrow{AB}$ et $\overrightarrow{AC}$.

On rappelle que, si $M(x_M\ ;\ y_M)$ et $N(x_N\ ;\ y_N)$, alors :
\[
\overrightarrow{MN}(x_N-x_M\ ;\ y_N-y_M).
\]

Ainsi :
\begin{align*}
\overrightarrow{AB}
&=(x_B-x_A\ ;\ y_B-y_A)\\
&=(3-(-1)\ ;\ 5-5)\\
&=(4\ ;\ 0).
\end{align*}

De même :
\begin{align*}
\overrightarrow{AC}
&=(x_C-x_A\ ;\ y_C-y_A)\\
&=(4-(-1)\ ;\ 0-5)\\
&=(5\ ;\ -5).
\end{align*}

Donc :
\[
\overrightarrow{AB}(4\ ;\ 0)
\qquad \text{et} \qquad
\overrightarrow{AC}(5\ ;\ -5).
\]

\item On utilise l'expression du produit scalaire avec les coordonnées.

Si $\vec u(x\ ;\ y)$ et $\vec v(x'\ ;\ y')$, alors :
\[
\vec u\cdot \vec v=xx'+yy'.
\]

Donc :
\begin{align*}
\overrightarrow{AB}\cdot\overrightarrow{AC}
&=4\times 5+0\times (-5)\\
&=20+0\\
&=20.
\end{align*}

Ainsi :
\[
\overrightarrow{AB}\cdot\overrightarrow{AC}=20.
\]
\end{enumerate}

\item \begin{enumerate}
\item On calcule la longueur $AC$ à l'aide de la formule de distance dans un repère orthonormé.

On a :
\begin{align*}
AC
&=\sqrt{(x_C-x_A)^2+(y_C-y_A)^2}\\
&=\sqrt{(4-(-1))^2+(0-5)^2}\\
&=\sqrt{5^2+(-5)^2}\\
&=\sqrt{25+25}\\
&=\sqrt{50}\\
&=\sqrt{25\times 2}\\
&=5\sqrt{2}.
\end{align*}

Donc :
\[
AC=5\sqrt{2}.
\]

\item On admet que $AB=4$.

L'expression du produit scalaire en fonction de l'angle $\widehat{BAC}$ est :
\[
\overrightarrow{AB}\cdot\overrightarrow{AC}
=AB\times AC\times \cos\left(\widehat{BAC}\right).
\]

En remplaçant $AB$ par $4$ et $AC$ par $5\sqrt{2}$, on obtient :
\[
\overrightarrow{AB}\cdot\overrightarrow{AC}
=4\times 5\sqrt{2}\times \cos\left(\widehat{BAC}\right).
\]

C'est-à-dire :
\[
\overrightarrow{AB}\cdot\overrightarrow{AC}
=20\sqrt{2}\cos\left(\widehat{BAC}\right).
\]

\item D'après la question 1.b, on sait que :
\[
\overrightarrow{AB}\cdot\overrightarrow{AC}=20.
\]

Or, d'après la question précédente :
\[
\overrightarrow{AB}\cdot\overrightarrow{AC}
=20\sqrt{2}\cos\left(\widehat{BAC}\right).
\]

On obtient donc :
\begin{align*}
20&=20\sqrt{2}\cos\left(\widehat{BAC}\right)\\
\cos\left(\widehat{BAC}\right)&=\frac{20}{20\sqrt{2}}\\
\cos\left(\widehat{BAC}\right)&=\frac{1}{\sqrt{2}}\\
\cos\left(\widehat{BAC}\right)&=\frac{\sqrt{2}}{2}.
\end{align*}

L'angle $\widehat{BAC}$ est un angle géométrique ; sa mesure appartient donc à l'intervalle $[0\ ;\ \pi]$.

Comme :
\[
\cos\left(\frac{\pi}{4}\right)=\frac{\sqrt{2}}{2},
\]
on en déduit que :
\[
\widehat{BAC}=\frac{\pi}{4}.
\]

Une mesure, en radian, de l'angle $\widehat{BAC}$ est donc :
\[
\boxed{\frac{\pi}{4}}.
\]
\end{enumerate}
\end{enumerate}

\clearpage
